% --- Archon physics formalization source begin ---
% archon:physics
% archon:covers PhyXMiniProblems/problem_phyx_mini_0827.lean
% archon:source-report reports/phyx_mini/problem_phyx_mini_0827.source.json

\chapter{Physics problem phyx\_mini\_0827}
\label{ch:PhyXMiniProblems_problem_phyx_mini_0827}

\paragraph{Problem source.}
\#\# Physical scenario

A door 1.00 m wide, of mass 15 kg, can rotate freely about a vertical axis through its hinges. A bullet with a mass of 10 g and a speed of 400 m/s strikes the center of the door as shown in figure, in a direction perpendicular to the plane of the door, and embeds itself there.

\#\# Question

Find the door's angular speed.

\#\# Answer choices

- A: A: \textbackslash{}( 0.60 \textbackslash{}text\{ rad/s\} \textbackslash{})

- B: B: \textbackslash{}( 0.40 \textbackslash{}text\{ rad/s\} \textbackslash{})

- C: C: \textbackslash{}( 0.80 \textbackslash{}text\{ rad/s\} \textbackslash{})

- D: D: \textbackslash{}( 0.30 \textbackslash{}text\{ rad/s\} \textbackslash{})

\#\# Figure caption (auxiliary; use the image as primary evidence)

The image depicts a physical scenario involving a bullet and a rod. Here's a detailed description:

- A vertical rod is shown with a hinge at the top, acting as a pivot point.
- The rod has a mass of \textbackslash{}(M = 15 \textbackslash{}, \textbackslash{}text\{kg\}\textbackslash{}) and a length labeled as \textbackslash{}(\textbackslash{}ell = 0.50 \textbackslash{}, \textbackslash{}text\{m\}\textbackslash{}).
- A bullet, labeled with a mass \textbackslash{}(m = 10 \textbackslash{}, \textbackslash{}text\{g\}\textbackslash{}), travels horizontally with a velocity \textbackslash{}(v\_\{\textbackslash{}text\{bullet\}\} = 400 \textbackslash{}, \textbackslash{}text\{m/s\}\textbackslash{}) toward the rod.
- The bullet is shown striking the rod at a point \textbackslash{}(d = 1.00 \textbackslash{}, \textbackslash{}text\{m\}\textbackslash{}) below the hinge.
- The “Before” and “After” segments indicate the positions of the rod and bullet before and after the collision, respectively.
- In the "After" state, the rod is at an angle to the right, with the bullet embedded in it and an arrow indicating rotation with angular velocity \textbackslash{}(\textbackslash{}omega\textbackslash{}).

The diagram visually conveys the dynamics of the collision and transfer of momentum that results in the rod's rotation.

\#\# Dataset metadata

Momentum and Collisions; Physical Model Grounding Reasoning, Multi-Formula Reasoning

\paragraph{Recorded answer/context.}
B: B: \textbackslash{}( 0.40 \textbackslash{}text\{ rad/s\} \textbackslash{})

\paragraph{Figure/image path.}
/root/proposal\_for\_physic/science-mango/phyx\_mini\_run/phyx\_data/test\_image/827.png

\paragraph{Formalization target.}
create a compiling Lean file with sorry bodies at `PhyXMiniProblems/problem\_phyx\_mini\_0827.lean`. The Lean declarations must preserve the physical quantities, dimensions or dimensional roles, figure labels, governing-law hypotheses, and final relation expressed by this problem.
Use Mathlib/Physlib names found through LeanExplore where available. If a physics API is missing, introduce faithful local abstractions rather than scalar placeholder aliases.

\begin{theorem}[Physics formalization target]
\label{thm:physics:phyx_mini_0827:target}
\lean{PhyXMiniProblems.ProblemPhyXMini0827.problem_phyx_mini_0827}
\uses{def:physics:phyx-mini-0827:phyxminiproblems-problemphyxmini0827-angularspeedinradianspersecond, def:physics:phyx-mini-0827:phyxminiproblems-problemphyxmini0827-doorbulletsetup, def:physics:phyx-mini-0827:phyxminiproblems-problemphyxmini0827-matchesdoorbulletscenario, def:physics:phyx-mini-0827:phyxminiproblems-problemphyxmini0827-matchesproblemreadouts, def:physics:phyx-mini-0827:phyxminiproblems-problemphyxmini0827-matchesprimarydoorimpactfigure, def:physics:phyx-mini-0827:phyxminiproblems-problemphyxmini0827-matchescollisionboundaryconditions, def:physics:phyx-mini-0827:phyxminiproblems-problemphyxmini0827-hasphysicaldoorbulletparameters, def:physics:phyx-mini-0827:phyxminiproblems-problemphyxmini0827-satisfiescenterimpactgeometry, def:physics:phyx-mini-0827:phyxminiproblems-problemphyxmini0827-satisfiesdoorbulletinertialaws, def:physics:phyx-mini-0827:phyxminiproblems-problemphyxmini0827-satisfiesangularmomentumconservationatimpact, def:physics:phyx-mini-0827:phyxminiproblems-problemphyxmini0827-recordeddatasetanswer, def:physics:phyx-mini-0827:phyxminiproblems-problemphyxmini0827-roundstonearesthundredthradianpersecond, def:physics:phyx-mini-0827:phyxminiproblems-problemphyxmini0827-matchesanswerchoice, def:physics:phyx-mini-0827:phyxminiproblems-problemphyxmini0827-isuniqueclosestdisplayedchoice}
The assigned autoformalize agent should translate this physics problem into Lean declarations in the covered file, with theorem and lemma proof bodies written as `by sorry`.
\end{theorem}
\begin{proof}
This is an autoformalization task, not a proof task. Produce statements that can later be proved by the physics prover without weakening the physical model.
\end{proof}

% --- Archon declaration topology begin ---
\section{Declaration topology}
\begin{definition}[moment Of Inertia Dimension]
  \label{def:physics:phyx-mini-0827:phyxminiproblems-problemphyxmini0827-momentofinertiadimension}
  \lean{PhyXMiniProblems.ProblemPhyXMini0827.momentOfInertiaDimension}
  The physical dimension mass * length\textasciicircum{}2 of an axial moment of inertia
\end{definition}
\begin{proof}
  This is the declared model definition of moment Of Inertia Dimension.
\end{proof}
\begin{definition}[angular Momentum Dimension]
  \label{def:physics:phyx-mini-0827:phyxminiproblems-problemphyxmini0827-angularmomentumdimension}
  \lean{PhyXMiniProblems.ProblemPhyXMini0827.angularMomentumDimension}
  \uses{def:physics:phyx-mini-0827:phyxminiproblems-problemphyxmini0827-momentofinertiadimension}
  The physical dimension mass * length\textasciicircum{}2 / time of angular momentum
\end{definition}
\begin{proof}
  This is the declared model definition of angular Momentum Dimension.
\end{proof}
\begin{definition}[Mass Quantity]
  \label{def:physics:phyx-mini-0827:phyxminiproblems-problemphyxmini0827-massquantity}
  \lean{PhyXMiniProblems.ProblemPhyXMini0827.MassQuantity}
  A nonnegative, unit-independent physical mass
\end{definition}
\begin{proof}
  This is the declared model definition of Mass Quantity.
\end{proof}
\begin{definition}[Length Quantity]
  \label{def:physics:phyx-mini-0827:phyxminiproblems-problemphyxmini0827-lengthquantity}
  \lean{PhyXMiniProblems.ProblemPhyXMini0827.LengthQuantity}
  A nonnegative, unit-independent physical length
\end{definition}
\begin{proof}
  This is the declared model definition of Length Quantity.
\end{proof}
\begin{definition}[Speed Quantity]
  \label{def:physics:phyx-mini-0827:phyxminiproblems-problemphyxmini0827-speedquantity}
  \lean{PhyXMiniProblems.ProblemPhyXMini0827.SpeedQuantity}
  A nonnegative, unit-independent speed magnitude
\end{definition}
\begin{proof}
  This is the declared model definition of Speed Quantity.
\end{proof}
\begin{definition}[Angular Speed Quantity]
  \label{def:physics:phyx-mini-0827:phyxminiproblems-problemphyxmini0827-angularspeedquantity}
  \lean{PhyXMiniProblems.ProblemPhyXMini0827.AngularSpeedQuantity}
  A nonnegative angular-speed magnitude; radians are dimensionless
\end{definition}
\begin{proof}
  This is the declared model definition of Angular Speed Quantity.
\end{proof}
\begin{definition}[Moment Of Inertia Quantity]
  \label{def:physics:phyx-mini-0827:phyxminiproblems-problemphyxmini0827-momentofinertiaquantity}
  \lean{PhyXMiniProblems.ProblemPhyXMini0827.MomentOfInertiaQuantity}
  \uses{def:physics:phyx-mini-0827:phyxminiproblems-problemphyxmini0827-momentofinertiadimension}
  A nonnegative scalar moment of inertia about the hinge axis
\end{definition}
\begin{proof}
  This is the declared model definition of Moment Of Inertia Quantity.
\end{proof}
\begin{definition}[Angular Momentum Quantity]
  \label{def:physics:phyx-mini-0827:phyxminiproblems-problemphyxmini0827-angularmomentumquantity}
  \lean{PhyXMiniProblems.ProblemPhyXMini0827.AngularMomentumQuantity}
  \uses{def:physics:phyx-mini-0827:phyxminiproblems-problemphyxmini0827-angularmomentumdimension}
  A nonnegative angular-momentum magnitude about the hinge axis
\end{definition}
\begin{proof}
  This is the declared model definition of Angular Momentum Quantity.
\end{proof}
\begin{definition}[mass Readout]
  \label{def:physics:phyx-mini-0827:phyxminiproblems-problemphyxmini0827-massreadout}
  \lean{PhyXMiniProblems.ProblemPhyXMini0827.massReadout}
  \uses{def:physics:phyx-mini-0827:phyxminiproblems-problemphyxmini0827-massquantity}
  Read a mass in a selected mass unit
\end{definition}
\begin{proof}
  This is the declared model definition of mass Readout.
\end{proof}
\begin{definition}[length Readout]
  \label{def:physics:phyx-mini-0827:phyxminiproblems-problemphyxmini0827-lengthreadout}
  \lean{PhyXMiniProblems.ProblemPhyXMini0827.lengthReadout}
  \uses{def:physics:phyx-mini-0827:phyxminiproblems-problemphyxmini0827-lengthquantity}
  Read a length in a selected length unit
\end{definition}
\begin{proof}
  This is the declared model definition of length Readout.
\end{proof}
\begin{definition}[speed Readout]
  \label{def:physics:phyx-mini-0827:phyxminiproblems-problemphyxmini0827-speedreadout}
  \lean{PhyXMiniProblems.ProblemPhyXMini0827.speedReadout}
  \uses{def:physics:phyx-mini-0827:phyxminiproblems-problemphyxmini0827-speedquantity}
  Read speed in coherent selected length and time units
\end{definition}
\begin{proof}
  This is the declared model definition of speed Readout.
\end{proof}
\begin{definition}[angular Speed Readout]
  \label{def:physics:phyx-mini-0827:phyxminiproblems-problemphyxmini0827-angularspeedreadout}
  \lean{PhyXMiniProblems.ProblemPhyXMini0827.angularSpeedReadout}
  \uses{def:physics:phyx-mini-0827:phyxminiproblems-problemphyxmini0827-angularspeedquantity}
  Read angular speed in inverse units of the selected time unit
\end{definition}
\begin{proof}
  This is the declared model definition of angular Speed Readout.
\end{proof}
\begin{definition}[moment Of Inertia Readout]
  \label{def:physics:phyx-mini-0827:phyxminiproblems-problemphyxmini0827-momentofinertiareadout}
  \lean{PhyXMiniProblems.ProblemPhyXMini0827.momentOfInertiaReadout}
  \uses{def:physics:phyx-mini-0827:phyxminiproblems-problemphyxmini0827-momentofinertiaquantity}
  Read moment of inertia in coherent selected mass and length units
\end{definition}
\begin{proof}
  This is the declared model definition of moment Of Inertia Readout.
\end{proof}
\begin{definition}[angular Momentum Readout]
  \label{def:physics:phyx-mini-0827:phyxminiproblems-problemphyxmini0827-angularmomentumreadout}
  \lean{PhyXMiniProblems.ProblemPhyXMini0827.angularMomentumReadout}
  \uses{def:physics:phyx-mini-0827:phyxminiproblems-problemphyxmini0827-angularmomentumquantity}
  Read angular momentum in coherent selected mechanical units
\end{definition}
\begin{proof}
  This is the declared model definition of angular Momentum Readout.
\end{proof}
\begin{definition}[mass In Kilograms]
  \label{def:physics:phyx-mini-0827:phyxminiproblems-problemphyxmini0827-massinkilograms}
  \lean{PhyXMiniProblems.ProblemPhyXMini0827.massInKilograms}
  \uses{def:physics:phyx-mini-0827:phyxminiproblems-problemphyxmini0827-massquantity, def:physics:phyx-mini-0827:phyxminiproblems-problemphyxmini0827-massreadout}
  Kilogram readout of a physical mass
\end{definition}
\begin{proof}
  This is the declared model definition of mass In Kilograms.
\end{proof}
\begin{definition}[mass In Grams]
  \label{def:physics:phyx-mini-0827:phyxminiproblems-problemphyxmini0827-massingrams}
  \lean{PhyXMiniProblems.ProblemPhyXMini0827.massInGrams}
  \uses{def:physics:phyx-mini-0827:phyxminiproblems-problemphyxmini0827-massquantity, def:physics:phyx-mini-0827:phyxminiproblems-problemphyxmini0827-massreadout}
  Gram readout used for the bullet-mass label
\end{definition}
\begin{proof}
  This is the declared model definition of mass In Grams.
\end{proof}
\begin{definition}[length In Meters]
  \label{def:physics:phyx-mini-0827:phyxminiproblems-problemphyxmini0827-lengthinmeters}
  \lean{PhyXMiniProblems.ProblemPhyXMini0827.lengthInMeters}
  \uses{def:physics:phyx-mini-0827:phyxminiproblems-problemphyxmini0827-lengthquantity, def:physics:phyx-mini-0827:phyxminiproblems-problemphyxmini0827-lengthreadout}
  Metre readout of a physical length
\end{definition}
\begin{proof}
  This is the declared model definition of length In Meters.
\end{proof}
\begin{definition}[speed In Meters Per Second]
  \label{def:physics:phyx-mini-0827:phyxminiproblems-problemphyxmini0827-speedinmeterspersecond}
  \lean{PhyXMiniProblems.ProblemPhyXMini0827.speedInMetersPerSecond}
  \uses{def:physics:phyx-mini-0827:phyxminiproblems-problemphyxmini0827-speedquantity, def:physics:phyx-mini-0827:phyxminiproblems-problemphyxmini0827-speedreadout}
  Metres-per-second readout of a speed magnitude
\end{definition}
\begin{proof}
  This is the declared model definition of speed In Meters Per Second.
\end{proof}
\begin{definition}[angular Speed In Radians Per Second]
  \label{def:physics:phyx-mini-0827:phyxminiproblems-problemphyxmini0827-angularspeedinradianspersecond}
  \lean{PhyXMiniProblems.ProblemPhyXMini0827.angularSpeedInRadiansPerSecond}
  \uses{def:physics:phyx-mini-0827:phyxminiproblems-problemphyxmini0827-angularspeedquantity, def:physics:phyx-mini-0827:phyxminiproblems-problemphyxmini0827-angularspeedreadout}
  Radians-per-second readout of an angular-speed magnitude
\end{definition}
\begin{proof}
  This is the declared model definition of angular Speed In Radians Per Second.
\end{proof}
\begin{definition}[moment Of Inertia In Kilogram Meters Squared]
  \label{def:physics:phyx-mini-0827:phyxminiproblems-problemphyxmini0827-momentofinertiainkilogrammeterssquared}
  \lean{PhyXMiniProblems.ProblemPhyXMini0827.momentOfInertiaInKilogramMetersSquared}
  \uses{def:physics:phyx-mini-0827:phyxminiproblems-problemphyxmini0827-momentofinertiaquantity, def:physics:phyx-mini-0827:phyxminiproblems-problemphyxmini0827-momentofinertiareadout}
  Kilogram-metre-squared readout of a moment of inertia
\end{definition}
\begin{proof}
  This is the declared model definition of moment Of Inertia In Kilogram Meters Squared.
\end{proof}
\begin{definition}[angular Momentum In Kilogram Meters Squared Per Second]
  \label{def:physics:phyx-mini-0827:phyxminiproblems-problemphyxmini0827-angularmomentuminkilogrammeterssquaredpersecond}
  \lean{PhyXMiniProblems.ProblemPhyXMini0827.angularMomentumInKilogramMetersSquaredPerSecond}
  \uses{def:physics:phyx-mini-0827:phyxminiproblems-problemphyxmini0827-angularmomentumquantity, def:physics:phyx-mini-0827:phyxminiproblems-problemphyxmini0827-angularmomentumreadout}
  Kilogram-metre-squared-per-second readout of angular momentum
\end{definition}
\begin{proof}
  This is the declared model definition of angular Momentum In Kilogram Meters Squared Per Second.
\end{proof}
\begin{definition}[Collision Stage]
  \label{def:physics:phyx-mini-0827:phyxminiproblems-problemphyxmini0827-collisionstage}
  \lean{PhyXMiniProblems.ProblemPhyXMini0827.CollisionStage}
  Configurations explicitly labelled Before and After in image 827.png
\end{definition}
\begin{proof}
  This is the declared model definition of Collision Stage.
\end{proof}
\begin{definition}[Figure Object]
  \label{def:physics:phyx-mini-0827:phyxminiproblems-problemphyxmini0827-figureobject}
  \lean{PhyXMiniProblems.ProblemPhyXMini0827.FigureObject}
  Physical or graphical objects visibly distinguished in the primary image
\end{definition}
\begin{proof}
  This is the declared model definition of Figure Object.
\end{proof}
\begin{definition}[Figure Label]
  \label{def:physics:phyx-mini-0827:phyxminiproblems-problemphyxmini0827-figurelabel}
  \lean{PhyXMiniProblems.ProblemPhyXMini0827.FigureLabel}
  Literal symbolic and stage labels visible in the primary image
\end{definition}
\begin{proof}
  This is the declared model definition of Figure Label.
\end{proof}
\begin{definition}[Door Rotation Axis]
  \label{def:physics:phyx-mini-0827:phyxminiproblems-problemphyxmini0827-doorrotationaxis}
  \lean{PhyXMiniProblems.ProblemPhyXMini0827.DoorRotationAxis}
  The physical axis about which the door is allowed to rotate
\end{definition}
\begin{proof}
  This is the declared model definition of Door Rotation Axis.
\end{proof}
\begin{definition}[Incidence Direction]
  \label{def:physics:phyx-mini-0827:phyxminiproblems-problemphyxmini0827-incidencedirection}
  \lean{PhyXMiniProblems.ProblemPhyXMini0827.IncidenceDirection}
  Direction of the incident bullet relative to the plane of the door
\end{definition}
\begin{proof}
  This is the declared model definition of Incidence Direction.
\end{proof}
\begin{definition}[Door Model]
  \label{def:physics:phyx-mini-0827:phyxminiproblems-problemphyxmini0827-doormodel}
  \lean{PhyXMiniProblems.ProblemPhyXMini0827.DoorModel}
  Idealized mass distribution used by the standard door-inertia law
\end{definition}
\begin{proof}
  This is the declared model definition of Door Model.
\end{proof}
\begin{definition}[Collision Outcome]
  \label{def:physics:phyx-mini-0827:phyxminiproblems-problemphyxmini0827-collisionoutcome}
  \lean{PhyXMiniProblems.ProblemPhyXMini0827.CollisionOutcome}
  Qualitative outcome of the short bullet-door collision
\end{definition}
\begin{proof}
  This is the declared model definition of Collision Outcome.
\end{proof}
\begin{definition}[Depicted Rotation Sense]
  \label{def:physics:phyx-mini-0827:phyxminiproblems-problemphyxmini0827-depictedrotationsense}
  \lean{PhyXMiniProblems.ProblemPhyXMini0827.DepictedRotationSense}
  Direction in which the after-impact door is drawn relative to Before
\end{definition}
\begin{proof}
  This is the declared model definition of Depicted Rotation Sense.
\end{proof}
\begin{definition}[Door Impact Figure]
  \label{def:physics:phyx-mini-0827:phyxminiproblems-problemphyxmini0827-doorimpactfigure}
  \lean{PhyXMiniProblems.ProblemPhyXMini0827.DoorImpactFigure}
  \uses{def:physics:phyx-mini-0827:phyxminiproblems-problemphyxmini0827-massquantity, def:physics:phyx-mini-0827:phyxminiproblems-problemphyxmini0827-lengthquantity, def:physics:phyx-mini-0827:phyxminiproblems-problemphyxmini0827-speedquantity, def:physics:phyx-mini-0827:phyxminiproblems-problemphyxmini0827-angularspeedquantity, def:physics:phyx-mini-0827:phyxminiproblems-problemphyxmini0827-collisionstage, def:physics:phyx-mini-0827:phyxminiproblems-problemphyxmini0827-figureobject, def:physics:phyx-mini-0827:phyxminiproblems-problemphyxmini0827-figurelabel, def:physics:phyx-mini-0827:phyxminiproblems-problemphyxmini0827-depictedrotationsense}
  Typed evidence transcribed from the primary bitmap. In particular, the independent fields for d and ℓ preserve the image's full-width and hinge-to-centre roles rather than following the swapped auxiliary caption
\end{definition}
\begin{proof}
  This is the declared model definition of Door Impact Figure.
\end{proof}
\begin{definition}[Door Bullet Setup]
  \label{def:physics:phyx-mini-0827:phyxminiproblems-problemphyxmini0827-doorbulletsetup}
  \lean{PhyXMiniProblems.ProblemPhyXMini0827.DoorBulletSetup}
  \uses{def:physics:phyx-mini-0827:phyxminiproblems-problemphyxmini0827-massquantity, def:physics:phyx-mini-0827:phyxminiproblems-problemphyxmini0827-lengthquantity, def:physics:phyx-mini-0827:phyxminiproblems-problemphyxmini0827-speedquantity, def:physics:phyx-mini-0827:phyxminiproblems-problemphyxmini0827-angularspeedquantity, def:physics:phyx-mini-0827:phyxminiproblems-problemphyxmini0827-momentofinertiaquantity, def:physics:phyx-mini-0827:phyxminiproblems-problemphyxmini0827-angularmomentumquantity, def:physics:phyx-mini-0827:phyxminiproblems-problemphyxmini0827-doorrotationaxis, def:physics:phyx-mini-0827:phyxminiproblems-problemphyxmini0827-incidencedirection, def:physics:phyx-mini-0827:phyxminiproblems-problemphyxmini0827-doormodel, def:physics:phyx-mini-0827:phyxminiproblems-problemphyxmini0827-collisionoutcome, def:physics:phyx-mini-0827:phyxminiproblems-problemphyxmini0827-doorimpactfigure}
  Independent physical quantities for the collision. The requested final angular speed is stored independently and is constrained only by the laws below; it is not defined to be an answer choice
\end{definition}
\begin{proof}
  This is the declared model definition of Door Bullet Setup.
\end{proof}
\begin{definition}[Matches Door Bullet Scenario]
  \label{def:physics:phyx-mini-0827:phyxminiproblems-problemphyxmini0827-matchesdoorbulletscenario}
  \lean{PhyXMiniProblems.ProblemPhyXMini0827.MatchesDoorBulletScenario}
  \uses{def:physics:phyx-mini-0827:phyxminiproblems-problemphyxmini0827-doorbulletsetup}
  Qualitative apparatus and collision assumptions stated in the problem
\end{definition}
\begin{proof}
  This is the declared model definition of Matches Door Bullet Scenario.
\end{proof}
\begin{definition}[Matches Problem Readouts]
  \label{def:physics:phyx-mini-0827:phyxminiproblems-problemphyxmini0827-matchesproblemreadouts}
  \lean{PhyXMiniProblems.ProblemPhyXMini0827.MatchesProblemReadouts}
  \uses{def:physics:phyx-mini-0827:phyxminiproblems-problemphyxmini0827-massinkilograms, def:physics:phyx-mini-0827:phyxminiproblems-problemphyxmini0827-massingrams, def:physics:phyx-mini-0827:phyxminiproblems-problemphyxmini0827-lengthinmeters, def:physics:phyx-mini-0827:phyxminiproblems-problemphyxmini0827-speedinmeterspersecond, def:physics:phyx-mini-0827:phyxminiproblems-problemphyxmini0827-doorbulletsetup}
  The values printed in the problem and primary image
\end{definition}
\begin{proof}
  This is the declared model definition of Matches Problem Readouts.
\end{proof}
\begin{definition}[Matches Primary Door Impact Figure]
  \label{def:physics:phyx-mini-0827:phyxminiproblems-problemphyxmini0827-matchesprimarydoorimpactfigure}
  \lean{PhyXMiniProblems.ProblemPhyXMini0827.MatchesPrimaryDoorImpactFigure}
  \uses{def:physics:phyx-mini-0827:phyxminiproblems-problemphyxmini0827-doorbulletsetup}
  Direct object, label, and spatial evidence from image 827.png
\end{definition}
\begin{proof}
  This is the declared model definition of Matches Primary Door Impact Figure.
\end{proof}
\begin{definition}[Matches Collision Boundary Conditions]
  \label{def:physics:phyx-mini-0827:phyxminiproblems-problemphyxmini0827-matchescollisionboundaryconditions}
  \lean{PhyXMiniProblems.ProblemPhyXMini0827.MatchesCollisionBoundaryConditions}
  \uses{def:physics:phyx-mini-0827:phyxminiproblems-problemphyxmini0827-angularspeedreadout, def:physics:phyx-mini-0827:phyxminiproblems-problemphyxmini0827-doorbulletsetup}
  Rest condition and hinge-impulse facts for the short collision interval
\end{definition}
\begin{proof}
  This is the declared model definition of Matches Collision Boundary Conditions.
\end{proof}
\begin{definition}[Has Physical Door Bullet Parameters]
  \label{def:physics:phyx-mini-0827:phyxminiproblems-problemphyxmini0827-hasphysicaldoorbulletparameters}
  \lean{PhyXMiniProblems.ProblemPhyXMini0827.HasPhysicalDoorBulletParameters}
  \uses{def:physics:phyx-mini-0827:phyxminiproblems-problemphyxmini0827-massinkilograms, def:physics:phyx-mini-0827:phyxminiproblems-problemphyxmini0827-lengthinmeters, def:physics:phyx-mini-0827:phyxminiproblems-problemphyxmini0827-speedinmeterspersecond, def:physics:phyx-mini-0827:phyxminiproblems-problemphyxmini0827-momentofinertiainkilogrammeterssquared, def:physics:phyx-mini-0827:phyxminiproblems-problemphyxmini0827-doorbulletsetup}
  Positivity selecting the nondegenerate physical branch
\end{definition}
\begin{proof}
  This is the declared model definition of Has Physical Door Bullet Parameters.
\end{proof}
\begin{definition}[Satisfies Center Impact Geometry]
  \label{def:physics:phyx-mini-0827:phyxminiproblems-problemphyxmini0827-satisfiescenterimpactgeometry}
  \lean{PhyXMiniProblems.ProblemPhyXMini0827.SatisfiesCenterImpactGeometry}
  \uses{def:physics:phyx-mini-0827:phyxminiproblems-problemphyxmini0827-lengthreadout, def:physics:phyx-mini-0827:phyxminiproblems-problemphyxmini0827-doorbulletsetup}
  The bullet strikes the door's centre, so its lever arm about the hinge axis is half the door width. This is a geometric law and does not mention omega
\end{definition}
\begin{proof}
  This is the declared model definition of Satisfies Center Impact Geometry.
\end{proof}
\begin{definition}[Satisfies Door Bullet Inertia Laws]
  \label{def:physics:phyx-mini-0827:phyxminiproblems-problemphyxmini0827-satisfiesdoorbulletinertialaws}
  \lean{PhyXMiniProblems.ProblemPhyXMini0827.SatisfiesDoorBulletInertiaLaws}
  \uses{def:physics:phyx-mini-0827:phyxminiproblems-problemphyxmini0827-massreadout, def:physics:phyx-mini-0827:phyxminiproblems-problemphyxmini0827-lengthreadout, def:physics:phyx-mini-0827:phyxminiproblems-problemphyxmini0827-momentofinertiareadout, def:physics:phyx-mini-0827:phyxminiproblems-problemphyxmini0827-doorbulletsetup}
  The uniform door has edge-axis inertia (1/3) M d\textasciicircum{}2; the embedded bullet is a point mass with inertia m ℓ\textasciicircum{}2; the after-impact inertia is their sum. These shape laws leave the requested angular speed unconstrained
\end{definition}
\begin{proof}
  This is the declared model definition of Satisfies Door Bullet Inertia Laws.
\end{proof}
\begin{definition}[Satisfies Angular Momentum Conservation At Impact]
  \label{def:physics:phyx-mini-0827:phyxminiproblems-problemphyxmini0827-satisfiesangularmomentumconservationatimpact}
  \lean{PhyXMiniProblems.ProblemPhyXMini0827.SatisfiesAngularMomentumConservationAtImpact}
  \uses{def:physics:phyx-mini-0827:phyxminiproblems-problemphyxmini0827-massreadout, def:physics:phyx-mini-0827:phyxminiproblems-problemphyxmini0827-lengthreadout, def:physics:phyx-mini-0827:phyxminiproblems-problemphyxmini0827-speedreadout, def:physics:phyx-mini-0827:phyxminiproblems-problemphyxmini0827-angularspeedreadout, def:physics:phyx-mini-0827:phyxminiproblems-problemphyxmini0827-momentofinertiareadout, def:physics:phyx-mini-0827:phyxminiproblems-problemphyxmini0827-angularmomentumreadout, def:physics:phyx-mini-0827:phyxminiproblems-problemphyxmini0827-doorbulletsetup}
  Angular momentum about the hinge axis during the short, perfectly inelastic impact. Perpendicular incidence makes the bullet term m v ℓ; the initially stationary door term is retained in the general balance and separately set to zero by the boundary conditions. The final embedded assembly has angular momentum I\_total * omega
\end{definition}
\begin{proof}
  This is the declared model definition of Satisfies Angular Momentum Conservation At Impact.
\end{proof}
\begin{lemma}[final Angular Speed exact readout]
  \label{lem:physics:phyx-mini-0827:phyxminiproblems-problemphyxmini0827-finalangularspeed-exact-readout}
  \lean{PhyXMiniProblems.ProblemPhyXMini0827.finalAngularSpeed_exact_readout}
  \uses{def:physics:phyx-mini-0827:phyxminiproblems-problemphyxmini0827-angularspeedinradianspersecond, def:physics:phyx-mini-0827:phyxminiproblems-problemphyxmini0827-doorbulletsetup, def:physics:phyx-mini-0827:phyxminiproblems-problemphyxmini0827-matchesproblemreadouts, def:physics:phyx-mini-0827:phyxminiproblems-problemphyxmini0827-matchescollisionboundaryconditions, def:physics:phyx-mini-0827:phyxminiproblems-problemphyxmini0827-hasphysicaldoorbulletparameters, def:physics:phyx-mini-0827:phyxminiproblems-problemphyxmini0827-satisfiescenterimpactgeometry, def:physics:phyx-mini-0827:phyxminiproblems-problemphyxmini0827-satisfiesdoorbulletinertialaws, def:physics:phyx-mini-0827:phyxminiproblems-problemphyxmini0827-satisfiesangularmomentumconservationatimpact}
  Substituting the printed values into the general inertia and conservation laws yields omega = (0.01 * 400 * 0.50) / ((1/3) * 15 * 1.00\textasciicircum{}2 + 0.01 * 0.50\textasciicircum{}2) = 800 / 2001. This exact angular speed is a derived conclusion, not a premise field
\end{lemma}
\begin{proof}
  The conclusion follows from the governing relations and figure readouts named in the statement of final Angular Speed exact readout.
\end{proof}
\begin{definition}[Answer Choice]
  \label{def:physics:phyx-mini-0827:phyxminiproblems-problemphyxmini0827-answerchoice}
  \lean{PhyXMiniProblems.ProblemPhyXMini0827.AnswerChoice}
  Labels of the four angular-speed choices printed in the source
\end{definition}
\begin{proof}
  This is the declared model definition of Answer Choice.
\end{proof}
\begin{definition}[displayed Angular Speed In Radians Per Second]
  \label{def:physics:phyx-mini-0827:phyxminiproblems-problemphyxmini0827-displayedangularspeedinradianspersecond}
  \lean{PhyXMiniProblems.ProblemPhyXMini0827.displayedAngularSpeedInRadiansPerSecond}
  \uses{def:physics:phyx-mini-0827:phyxminiproblems-problemphyxmini0827-answerchoice}
  Numerical value beside each answer label, in radians per second
\end{definition}
\begin{proof}
  This is the declared model definition of displayed Angular Speed In Radians Per Second.
\end{proof}
\begin{definition}[recorded Dataset Answer]
  \label{def:physics:phyx-mini-0827:phyxminiproblems-problemphyxmini0827-recordeddatasetanswer}
  \lean{PhyXMiniProblems.ProblemPhyXMini0827.recordedDatasetAnswer}
  \uses{def:physics:phyx-mini-0827:phyxminiproblems-problemphyxmini0827-answerchoice}
  The answer label recorded by the source dataset; this is metadata only
\end{definition}
\begin{proof}
  This is the declared model definition of recorded Dataset Answer.
\end{proof}
\begin{definition}[Rounds To Nearest Hundredth Radian Per Second]
  \label{def:physics:phyx-mini-0827:phyxminiproblems-problemphyxmini0827-roundstonearesthundredthradianpersecond}
  \lean{PhyXMiniProblems.ProblemPhyXMini0827.RoundsToNearestHundredthRadianPerSecond}
  \uses{def:physics:phyx-mini-0827:phyxminiproblems-problemphyxmini0827-angularspeedquantity, def:physics:phyx-mini-0827:phyxminiproblems-problemphyxmini0827-angularspeedinradianspersecond}
  Agreement with a value displayed to the nearest hundredth of a radian per second. Half a hundredth is 1/200 rad/s
\end{definition}
\begin{proof}
  This is the declared model definition of Rounds To Nearest Hundredth Radian Per Second.
\end{proof}
\begin{definition}[Matches Answer Choice]
  \label{def:physics:phyx-mini-0827:phyxminiproblems-problemphyxmini0827-matchesanswerchoice}
  \lean{PhyXMiniProblems.ProblemPhyXMini0827.MatchesAnswerChoice}
  \uses{def:physics:phyx-mini-0827:phyxminiproblems-problemphyxmini0827-doorbulletsetup, def:physics:phyx-mini-0827:phyxminiproblems-problemphyxmini0827-answerchoice, def:physics:phyx-mini-0827:phyxminiproblems-problemphyxmini0827-displayedangularspeedinradianspersecond, def:physics:phyx-mini-0827:phyxminiproblems-problemphyxmini0827-roundstonearesthundredthradianpersecond}
  The physical final angular speed rounds to the value printed by a choice
\end{definition}
\begin{proof}
  This is the declared model definition of Matches Answer Choice.
\end{proof}
\begin{definition}[Is Unique Closest Displayed Choice]
  \label{def:physics:phyx-mini-0827:phyxminiproblems-problemphyxmini0827-isuniqueclosestdisplayedchoice}
  \lean{PhyXMiniProblems.ProblemPhyXMini0827.IsUniqueClosestDisplayedChoice}
  \uses{def:physics:phyx-mini-0827:phyxminiproblems-problemphyxmini0827-angularspeedinradianspersecond, def:physics:phyx-mini-0827:phyxminiproblems-problemphyxmini0827-doorbulletsetup, def:physics:phyx-mini-0827:phyxminiproblems-problemphyxmini0827-answerchoice, def:physics:phyx-mini-0827:phyxminiproblems-problemphyxmini0827-displayedangularspeedinradianspersecond}
  A displayed choice is strictly closer than each of the alternatives
\end{definition}
\begin{proof}
  This is the declared model definition of Is Unique Closest Displayed Choice.
\end{proof}
% --- Archon declaration topology end ---
% --- Archon physics formalization source end ---
